\chapter{Comprehensive Metabolic Panel (CMP)}

\section*{What Is This Test?}
A comprehensive metabolic panel (CMP) is a group of 14 tests performed on a single blood sample that provides important information about your body's chemical balance and metabolism. It includes: glucose (blood sugar), calcium, total protein, albumin, bilirubin, alkaline phosphatase, AST, ALT, sodium, potassium, chloride, carbon dioxide (bicarbonate), blood urea nitrogen (BUN), and creatinine. The CMP evaluates your blood sugar, kidney function, liver function, and electrolyte and acid-base balance. It is one of the most commonly ordered laboratory tests and is used as a general health screening tool, to monitor chronic conditions such as diabetes and kidney disease, and to evaluate symptoms such as fatigue, nausea, or confusion \cite{wallach2022interpretation}.

\section*{Alternative Names}
CMP; Comprehensive Metabolic Panel; Chemistry Panel; Chem 14; SMA-14

\section*{Principle}
The CMP is performed on automated clinical chemistry analyzers using photometric (colorimetric) and potentiometric methods. Glucose is measured by enzymatic methods (glucose oxidase or hexokinase). Electrolytes (sodium, potassium, chloride) are measured by ion-selective electrode potentiometry. Kidney function markers (BUN, creatinine) are measured enzymatically. Liver function tests (ALT, AST, ALP, bilirubin) use kinetic photometric methods. Total protein and albumin are measured by dye-binding (bromocresol green for albumin, biuret for total protein). Calcium is measured by colorimetric methods or ion-selective electrodes. Carbon dioxide is measured by enzymatic or conductimetric methods. Results are typically available within 1 to 2 hours on modern high-throughput analyzers such as the Roche cobas c702, Siemens Atellica, or Abbott Alinity c \cite{tietz2023textbook}.

\section*{History}
The multi-test chemistry panel evolved from the AutoAnalyzer, developed by Leonard Skeggs at the Veterans Administration in the 1950s. The first automated chemistry panel could measure multiple analytes from a single blood sample. The SMA (Sequential Multiple Analyzer) series introduced in the 1960s-1970s by Technicon enabled sequential measurement of 12-20 chemistry tests. The SMA-12 (12 analytes) became a standard screening test. The term "comprehensive metabolic panel" was standardized by the College of American Pathologists (CAP) in the 1990s as a 14-test panel. Modern chemistry analyzers can perform over 400 tests per hour from microvolume samples. Point-of-care chemistry testing is now available for some analytes (glucose, electrolytes) at the bedside \cite{clsi2023laboratory}.

\section*{Why This Test Is Ordered}
\begin{itemize}
  \item Routine health screening: the CMP is one of the most commonly ordered laboratory tests, used as a general health assessment
  \item Evaluation of symptoms suggesting metabolic, kidney, or liver disease: fatigue, nausea, vomiting, jaundice, confusion, weight loss, or abdominal pain
  \item Assessment and monitoring of kidney function in patients with hypertension, diabetes, heart failure, or chronic kidney disease
  \item Evaluation of liver function and screening for liver disease: abnormal liver enzymes, hepatitis, fatty liver disease, or alcohol use disorder
  \item Assessment of glucose control and evaluation of hyperglycemia or hypoglycemia
  \item Preoperative assessment and hospital admission workup
  \item Monitoring of medications that affect electrolytes, kidney function, or the liver: diuretics, ACE inhibitors, ARBs, statins, anticonvulsants, and chemotherapy
  \item Monitoring of chronic conditions including diabetes, cirrhosis, and chronic kidney disease
  \item Evaluation of dehydration, acid-base disturbances, and electrolyte abnormalities in the emergency setting
  \item Screening for chronic kidney disease in at-risk populations
\end{itemize}

\section*{Primary vs Secondary Test}
This is a primary screening test for metabolic, kidney, and liver disorders.

\section*{How This Test Is Performed}
For most patients, a health care professional will take a blood sample from a vein in your arm using a small needle. After the needle is inserted, a small amount of blood will be collected into a test tube or vial. You may feel a little sting when the needle goes in or out. This usually takes less than five minutes. For certain tests, a urine sample, stool sample, or other specimen may be required.

\section*{What Preparation Is Needed}
Preparation requirements vary by test. Some tests require fasting for 8 to 12 hours. Others have no special preparation. Your healthcare provider or laboratory will inform you of any specific preparation needed for this test.

\section*{Contraindications and Precautions}
No absolute contraindications. Factors affecting results: Fasting status affects glucose (8-12 hour fast recommended for accurate glucose). Meal timing affects glucose, triglycerides, and electrolytes. Dehydration elevates BUN, creatinine, and electrolytes. Vigorous exercise may elevate AST, CK, and potassium. Specimen hemolysis falsely elevates potassium and LDH. Lipemia may interfere with some photometric assays. Medications may affect various analytes.

\section*{Risks}
Minimal risk. Slight pain or bruising at the venipuncture site.

\section*{What Happens After the Test}
Results are typically available within 1 to 2 hours. Abnormal results may prompt: Repeat testing to confirm findings, additional targeted tests (hepatitis serology if liver enzymes elevated, A1c if glucose elevated, lipid panel, urinalysis), referral to specialists (nephrology for kidney disease, endocrinology for diabetes, gastroenterology for liver disease). Your healthcare provider will interpret results in the context of your symptoms, medical history, and medications \cite{tierney2023current}.

\section*{Understanding Results}

\subsection*{Below Normal}
\begin{itemize}
  \item Low glucose: May indicate fasting, insulinoma, liver disease, or sepsis.
  \item Low albumin: May indicate malnutrition, liver disease, nephrotic syndrome, or chronic inflammation.
  \item Low calcium: May indicate hypoparathyroidism, vitamin D deficiency, or hypoalbuminemia (corrected calcium should be calculated).
  \item Low sodium: May indicate dilutional states (heart failure, SIADH) or sodium loss.
  \item Low potassium: May indicate GI losses, diuretic use, or metabolic alkalosis.
  \item Low bicarbonate: May indicate metabolic acidosis (diabetic ketoacidosis, lactic acidosis, renal tubular acidosis).
\end{itemize}

\subsection*{Above Normal}
\begin{itemize}
  \item Elevated glucose: May indicate diabetes mellitus, stress hyperglycemia, corticosteroid use, or hyperthyroidism.
  \item Fasting glucose >126 mg/dL on two occasions indicates diabetes.
  \item Elevated BUN and creatinine: Indicate impaired kidney function; the BUN/creatinine ratio helps distinguish causes (e.g., ratio >20:1 suggests prerenal azotemia such as dehydration or heart failure) \cite{fischbach2023manual}.
  \item Elevated liver enzymes (ALT, AST): May indicate liver injury from viral hepatitis, alcohol, medications, fatty liver disease, or biliary obstruction.
  \item Elevated ALP: May indicate biliary obstruction or bone disease.
  \item Elevated bilirubin: May indicate liver disease, hemolysis, or biliary obstruction.
  \item Electrolyte abnormalities: Hyponatremia (low sodium) may indicate SIADH, heart failure, or kidney disease; hyperkalemia may indicate kidney failure or cell lysis; hypernatremia indicates dehydration.
\end{itemize}

\section*{Normal Results}
\begin{itemize}
  \item \textbf{Glucose (fasting):} 70--100 mg/dL
  \item \textbf{BUN:} 7--20 mg/dL
  \item \textbf{Creatinine:} Males: 0.7--1.3 mg/dL; Females: 0.6--1.1 mg/dL
  \item \textbf{Sodium:} 136--145 mEq/L
  \item \textbf{Potassium:} 3.5--5.0 mEq/L
  \item \textbf{Chloride:} 98--106 mEq/L
  \item \textbf{Carbon dioxide:} 23--29 mEq/L
  \item \textbf{Calcium:} 8.5--10.5 mg/dL
  \item \textbf{Total protein:} 6.0--8.3 g/dL
  \item \textbf{Albumin:} 3.4--5.4 g/dL
  \item \textbf{Bilirubin (total):} 0.1--1.2 mg/dL
  \item \textbf{ALP:} 44--147 IU/L
  \item \textbf{AST:} 10--40 IU/L
  \item \textbf{ALT:} 7--56 IU/L
\end{itemize}

\section*{Follow-up Tests}
\begin{itemize}
  \item \textbf{Hemoglobin A1c:} If glucose is elevated, to assess long-term glycemic control
  \item \textbf{Lipid panel:} Often ordered alongside CMP for cardiovascular risk assessment
  \item \textbf{Urinalysis:} If kidney function is abnormal, to assess for proteinuria, hematuria, or casts
  \item \textbf{Hepatitis serology:} If liver enzymes are elevated, to evaluate for viral hepatitis
  \item \textbf{TSH:} If electrolyte or liver enzyme abnormalities are unexplained
  \item \textbf{Estimated GFR (eGFR):} Calculated from creatinine, age, and sex --- provides better assessment of kidney function than creatinine alone
\end{itemize}


\section*{References}
\bibliographystyle{unsrtnat}
\bibliography{references}

\clearpage
\section*{Regulatory Status and Authorization Date}
Regulatory status applies to the specific assay, instrument, manufacturer, and intended use rather than to the test name alone. No single approval date can be assigned to this test category. For a commercial assay, verify the current product in the relevant regulator's database: the U.S. Food and Drug Administration (FDA) clearance, approval, or authorization; India's Central Drugs Standard Control Organisation (CDSCO) licence or permission; the European Union In Vitro Diagnostic Regulation (EU IVDR) CE certificate issued through the applicable conformity-assessment route; or the United Kingdom Medicines and Healthcare products Regulatory Agency (MHRA) registration and UKCA/CE status. Laboratory-developed tests may not have an individual FDA, CDSCO, EU IVDR, or MHRA approval date.
\clearpage
